\section{Method}
\label{sec:method}

We formulate calibration as a trajectory-conditioned inverse problem.  Let
$z_t$ denote the simulated particle state, $u_t$ the measured poses and
velocities of the two spatulas, $\theta$ the fitted material parameters, and
$\psi$ the fixed reconstruction, contact, camera, and numerical settings:
\begin{equation}
 z_{t+1}=\mathcal{S}_{h,\Delta t}(z_t,u_t;\theta,\psi),\qquad
 y_t=\mathcal{H}_{c}(z_t)+\delta_t+\epsilon_t .
 \label{eq:inverse_problem}
\end{equation}
Here $\mathcal{H}_{c}$ projects particles into the calibrated depth camera,
while $\delta_t$ and $\epsilon_t$ denote model discrepancy and measurement
noise.  The estimated values are therefore effective parameters of this
specified experiment and simulator, rather than universal dough constants.

\subsection{Metric Reconstruction and Tool Replay}

A single RGB-D camera records depth and colour, while the motion-capture
system records the two spatula frames.  Camera intrinsics and a rigid
camera-to-scene transform are obtained from repeated AprilTag/mocap
correspondences.  The calibration also stores the table plane.  All point
clouds are transformed once into metric scene coordinates; no scale,
translation, or rigid registration is fitted during evaluation.

The first observed dough surface is voxelised and each occupied vertical
column is filled down to the calibrated table.  Given reconstructed volume
$V$, fixed density $\rho$, and $N$ sampled material points, the dataset uses
\begin{equation}
 M=\rho V,\qquad V_p=V/N,\qquad m_p=M/N .
 \label{eq:mass_initialization}
\end{equation}
Thus, episode mass is derived from reconstructed volume, and total volume and
mass remain unchanged when particle count is varied.  Tool positions are interpolated linearly and quaternion
orientations by shortest-path spherical interpolation.  The UR5e and Gen3
spatulas are represented by signed-distance fields generated from registered
meshes.  Their velocity at a contact point $x$ is
\begin{equation}
 v_{\mathrm{tool}}(x)=v_{\mathrm{lin}}+\omega\times(x-c),
 \label{eq:tool_velocity}
\end{equation}
where $c$ is the tool centre.  The unilateral Coulomb approximation removes
inward normal relative velocity and limits tangential correction using the
configured friction coefficient; separating motion is left unchanged.  A
post-advection projection resolves residual floor or tool penetration.

\subsection{Differentiable MLS-MPM}

The simulator implements three-dimensional MLS-MPM with APIC transfer
\cite{jiang2016mpm} in Taichi \cite{hu2019taichi}.  Each particle stores
position $x_p$, velocity $v_p$, deformation gradient $F_p$, affine velocity
matrix $C_p$, and optional plastic-volume state.  The deformation update is
\begin{equation}
 F_p^{n+1}=(I+\Delta t\,C_p)F_p^n .
 \label{eq:deformation_update}
\end{equation}
Young's modulus $E$ and Poisson ratio $\nu$ define
\begin{equation}
 \mu=\frac{E}{2(1+\nu)},\qquad
 \lambda=\frac{E\nu}{(1+\nu)(1-2\nu)} .
 \label{eq:lame}
\end{equation}
With polar factor $R$, $J=\det(F_p)$, and viscosity $\eta$, the implemented
Kirchhoff-style stress is
\begin{equation}
 \tau=2\mu(F_p-R)F_p^T+\lambda J(J-1)I
   +\eta(C_p+C_p^T).
 \label{eq:stress}
\end{equation}
The final term is an additive rate stress, not a multi-time-scale rheological
model.  For quadratic B-spline weight $w_{pi}$, particle-to-node offset
$d_{pi}$, and grid spacing $\Delta x$, particle-to-grid transfer accumulates
\begin{align}
 m_i &\mathrel{+}=w_{pi}m_p,\\
 m_i v_i &\mathrel{+}=w_{pi}\!\left[m_pv_p+
 \left(-4\Delta t\,V_p\Delta x^{-2}\tau+m_pC_p\right)d_{pi}\right].
 \label{eq:p2g}
\end{align}
The corrected grid-to-particle update reconstructs the affine velocity using
$4\Delta x^{-2}d_{pi}$.  Optional plasticity decomposes
$F_p=U\Sigma V^T$ and clamps each principal stretch to
$[\sigma_{\min},\sigma_{\max}]$ before reconstruction.  These bounds are
dimensionless stretch limits, not yield stresses.

\subsection{Partial-View Objective and Optimization}

At each selected frame, a differentiable camera-space particle renderer uses
compact splats and soft visibility to produce depth, coverage, and visible
particle coordinates in the calibrated camera.  After normalizing the
configured weights by their sum, the per-frame objective is
\begin{equation}
 \mathcal{L}_t=0.5\ell_d+0.25\ell_c+0.25\ell_q ,
 \label{eq:training_loss}
\end{equation}
where $\ell_d$ compares simulated and observed depth change relative to the
initial frame, $\ell_c$ penalizes missing positive observed foreground, and
$\ell_q$ is the directed distance from observed points to visible predicted
particles.  Unknown pixels are not treated as empty space.  Depth and point
distances are divided by the declared $0.01\,$m scales before applying the
Huber penalty.  This balances objective components; it does not normalize
physical material units.  Splat support, visibility, and nearest-particle
associations are held fixed within each local derivative evaluation.

Reverse-mode differentiation propagates adjoints through position, velocity,
affine velocity, deformation, and plastic history.  Host checkpoints bound
device memory; each segment recomputes forward intermediates and must normally
reproduce endpoint states and contact decisions.  Physical gradients are
mapped to optimizer coordinates: $E$ uses $q_E=\log(E/E_0)$, while $\nu$,
$\eta$, and the two stretch limits use bounded scaled coordinates.  Projected
Adam proposes updates within physical bounds, and Armijo backtracking accepts
only finite sufficient-decrease steps.  This coordinate transformation avoids
comparing raw gradients expressed in Pa, Pa\,s, and dimensionless units.

We also report an existing derivative-free baseline that varies only $E$ on
a logarithmic coarse-to-fine grid while holding its own trajectory setup and
other configured quantities fixed.  Because it uses different fixed settings,
it is evidence for the static search procedure rather than a matched optimizer
comparison. Candidate selection uses training frames only. After selection, a separate frozen
full-resolution evaluator reports depth, silhouette, coverage, and directed
point-distance metrics on the training interval and subsequent temporal
continuation.  Its strict metric is intentionally distinct from
Eq.~\eqref{eq:training_loss}; numerical values from the two objectives are not
interchangeable.

Finally, the interpretation separates four parameter families: material
$(E,\nu,\eta,\sigma_{\min},\sigma_{\max})$; interface quantities such as tool
friction and floor response; sensing/reconstruction quantities including
camera pose, hidden volume, mass, and tool registration; and numerical
quantities $(N,h,\Delta t)$ and SDF resolution.  The present estimates are
conditional on all fixed quantities in the latter three families.
